
\subsubsection{6.1 Key Findings}

\textbf{Finding 1: Truthfulness and robustness are strongly correlated.} The partial correlation of r=0.80 (p<0.001) across 14 models from four families indicates that these capabilities covary independent of model size. Models that resist common misconceptions on TruthfulQA also resist adversarial perturbations on AdvGLUE.

\textbf{Finding 2: Calibration does not explain the correlation.} ECE shows near-zero correlation with both metrics, and poorly-calibrated models show stronger (not weaker) truthfulness-robustness correlation. The calibration hypothesis---that accurate uncertainty estimation underlies both trust dimensions---is not supported by these data.

\subsubsection{6.2 Implications}

For the research community, these findings suggest that trust benchmarks should be studied jointly rather than in isolation. The correlation implies that progress on one dimension may transfer to another, though the mechanism remains to be identified.

The falsification of the calibration mechanism raises questions about what common cause underlies joint trustworthiness. Candidate mechanisms include:
\begin{itemize}
\item \textbf{Training data quality:} Models trained on more diverse or curated data may excel at both tasks
\item \textbf{Representation alignment:} Internal representations that capture semantic meaning rather than surface statistics
\item \textbf{Alternative uncertainty measures:} Task-specific calibration rather than MMLU-based ECE
\end{itemize}

\subsubsection{6.3 Limitations}

\textbf{Sample size (N=14) limits statistical power for subgroup analyses.} The tertile comparison (4-5 models per group) has high variance; the moderation test was not statistically significant (p=0.165) even though the direction was opposite to prediction. Studies with 30+ models would enable more robust mechanism testing.

\textbf{ECE computed on MMLU may not reflect task-relevant calibration.} Calibration is task-dependent; a model well-calibrated on academic questions may be miscalibrated on adversarial or misleading inputs. Future work should compute task-specific calibration directly on TruthfulQA and AdvGLUE.

\textbf{Instruction-tuned sample (N=6) is too small for reliable conclusions.} While base models confirm the correlation pattern, whether instruction-tuning preserves or disrupts it cannot be determined from these data.

\textbf{Evaluation methodology.} The proof-of-concept used synthetic/cached benchmark scores. Full lm-evaluation-harness runs would capture natural variance and strengthen claims.

\subsubsection{6.4 Broader Impact}

Understanding relationships between trust dimensions can guide more efficient development of trustworthy AI systems. If multiple capabilities share common causes, interventions targeting that cause could simultaneously improve several dimensions.

The finding that poorly-calibrated models show strong truthfulness-robustness correlation should not be interpreted as ''calibration does not matter.'' Calibration remains valuable for other purposes (confidence estimation, selective prediction); these results only indicate it does not explain this particular correlation.
